\section{最大安全接入容量计算与逆变器协调算法}\label{sec:algorithm}

本文的设计变量包括新增负荷容量 $L$ 和光伏逆变器的下垂增益向量 $\bm g=(g_1,\ldots,g_m)$。针对一组给定增益，算法首先计算馈线在当前负荷水平下的非线性 DistFlow；再利用分块 Hessian、半圆支撑函数及复功率几何界检查安全约束，并沿新增负荷方向搜索最大可接受容量。完成容量评估后，根据最紧约束所在节点及其灵敏度更新逆变器增益，重新运行安全容量计算。上述两个步骤交替执行，直到容量改善趋于平稳或达到计算预算。

\begin{figure}[htbp]
\centering
\begin{tikzpicture}[node distance=3.8mm,>=Latex,
 every node/.style={font=\small},
 process/.style={draw=blue!65!black,rounded corners=2pt,fill=blue!5,align=center,minimum height=9mm,text width=11cm,inner sep=3pt},
 decision/.style={draw=orange!80!black,rounded corners=2pt,fill=orange!7,align=center,minimum height=8mm,text width=11cm,inner sep=3pt}]
\node[process] (a) {输入馈线参数、光伏不确定集、逆变器增益和待接入容量 $L$};
\node[process,below=of a] (b) {按新增负荷划分单元 $I$；求解非线性 DistFlow，计算 Jacobian 与分块 Hessian};
\node[process,below=of b] (c) {半圆支撑函数计算电压边界；复功率角度法计算电流边界；检查完整 $P$--$Q$ 功率圆};
\node[decision,below=of c] (d) {当前完整负荷单元是否满足全部安全判据？};
\node[process,below=of d] (e) {满足：接受单元；未满足：细分并复算（达到预算则记为未知）};
\node[process,below=of e] (f) {记录可接入容量与活跃约束，更新逆变器增益后重新计算};
\foreach \s/\t in {a/b,b/c,c/d,d/e,e/f}{\draw[->,thick,blue!70!black] (\s.south)--(\t.north);}
\draw[->,thick,dashed,teal!70!black] (f.east) -- ++(.55,0) |- (b.east);
\end{tikzpicture}
\caption{最大安全接入容量与协调增益联合设计流程。每次增益更新后均重新完成负荷区间检查。}
\label{fig:flow}
\end{figure}


\subsection{完整负荷区间安全检查}
对于当前增益 $\bm g$，给定候选容量 $L$，把 $[0,L]$ 分成宽度不超过 $\Delta$ 的连续单元 $I_k=[a_k,b_k]$。每个单元完成以下三类计算：
\begin{enumerate}[leftmargin=8mm,itemsep=2pt]
\item 在代表运行点求解非线性 DistFlow，利用式~\eqref{eq:zt}--\eqref{eq:blocks} 计算电压对负荷与光伏扰动的 Jacobian 和分块 Hessian。
\item 利用式~\eqref{eq:voltage_accept} 计算节点电压裕度，利用式~\eqref{eq:phasor_bound} 计算线路电流裕度，并通过式~\eqref{eq:PVpass} 检查逆变器总视在功率是否落在运行范围内。
\item 当某个单元的安全上界不能满足要求时，先缩小单元宽度并重新计算。只有所有单元的全部约束均通过后，才接受候选容量 $L$。
\end{enumerate}
完整负荷区间检查比仅在若干独立容量点求解潮流更稳定：即使两个相邻采样点满足电压约束，中间区间仍可能出现更紧的约束或更大的非线性误差。

\subsection{协调增益与容量交替更新}
在本文算例中，以统一增益作为初始解，并允许四台逆变器在 $0\le g_j\le4$ 的研究范围内分别调整。采用坐标更新结合局部约束优化，每次提出一个新的增益组合后，都重新计算完整负荷前缀的安全容量。只保留容量改善且安全约束检查通过的方案；当各增益坐标的试探均不能继续提高目标值时，输出当前最佳设计。

\paragraph{算法 1：结构化安全容量与协调增益联合计算}
\begin{enumerate}[leftmargin=8mm,itemsep=2pt]
\item \textbf{输入：}网络参数、光伏不确定范围、设备约束、初始增益 $\bm g^{(0)}$、负荷划分宽度 $\Delta$ 和角度数 $N_\theta$。
\item \textbf{潮流及灵敏度：}对当前增益求解 DistFlow，并计算电压与线路复功率的线性响应和分块 Hessian 余项。
\item \textbf{容量搜索：}沿负荷方向逐步增大 $L$；对 $[0,L]$ 执行完整单元检查，针对临界单元自适应细分，记录首个难以通过的约束及当前可接受容量。
\item \textbf{协调更新：}识别活跃的电压、线路电流和逆变器视在功率约束；通过局部增益搜索形成新的 $\bm g$。
\item \textbf{重新评估：}将更新后的 $\bm g$ 代回步骤 2--3。若对应容量提高，则更新当前设计；否则调整搜索方向或步长。
\item \textbf{输出：}逆变器增益配置、接入容量、最紧约束、潮流求解记录以及完整负荷分段结果。
\end{enumerate}

本算法的容量内层和增益外层使用\emph{同一物理网络及同一安全约束}；增益不会根据某个孤立光伏场景的电压表现直接选定，而是由容量计算结果驱动。对于同样的增益设置，更紧的电流上界能够减少保守拒绝，进而增加可接受容量。

\subsection{计算结果与实施参数}

完成一轮容量搜索后，程序输出新增负荷容量 $L$、各台逆变器增益 $\bm g$、所有节点的电压上下限裕度、各条线路的电流裕度，以及逆变器视在功率裕度。每个候选容量还保存通过检查的负荷单元及求解次数，便于比较不同控制配置和计算参数。

算法的主要设置包括负荷单元宽度 $\Delta$、最小细分宽度和线路复功率角度数 $N_\theta$。较小的负荷单元能够提高临界容量附近的计算精度；较多的角度有助于收紧线路电流估计，同时增加计算量。第~\ref{sec:experiments} 节在 IEEE 33 节点算例中比较这些设置对安全接入容量及运行时间的影响。
